\section{How not to fool yourself}
\label{sec:pitfalls}

Interpretability has no ground truth. Nobody can open a language model and look up the right answer, so an
explanation can only be checked against its own predictions. The history of the field contains many results
that looked convincing and turned out to be artefacts. This section collects the main ways in which that
happens, each with the control experiment that guards against it. Most of them have already appeared in the
Caution boxes.

\subsection{Eight failure modes}

\paragraph{1. Decodable is not used.} A probe or a latent shows that information is present. It does not show
that the model reads it (\cref{sec:probing}). \emph{Control:} intervene on the representation and check that the
behaviour changes as predicted.

\paragraph{2. The method explains the data and not the model.} Some saliency maps survive the replacement of
the weights by random numbers \citep{adebayo2018sanity}, and SAEs find readable latents in randomly initialised
transformers \citep{heap2025sparse}. \emph{Control:} run the identical pipeline on a model with random weights
and on shuffled labels. Whatever survives is not about what the model has learned.

\paragraph{3. The strongest examples flatter.} The top activations of a neuron or latent can suggest a clean
concept that fails on the bulk of its activations and on other data sets
\citep{bolukbasi2021interpretability}. \emph{Control:} sample examples from the whole range of activations, use
a second corpus, and measure how often the concept occurs \emph{without} the unit firing.

\paragraph{4. The intervention leaves the training distribution.} Zero ablation and strong steering produce
activations that the model has never seen. The resulting behaviour may reflect the damage and not the role of
the component (\cref{sec:interventions}). \emph{Control:} prefer replacement values that occur naturally, report
the size of the intervention relative to typical activations, and check that unrelated outputs stay intact.

\paragraph{5. Compensation hides importance.} In \cref{tab:heads}, corrupting the most important head changed
the output by a fifth of its direct effect, because later layers made up for it. \emph{Control:} report direct
and total effects, and ablate groups of components that may substitute for one another.

\paragraph{6. A successful intervention through the wrong pathway.} A direction or subspace found by
optimisation can produce the predicted behaviour by activating something the model does not normally use
\citep{makelov2024subspace}. \emph{Control:} verify that the direction varies in normal operation in the way the
hypothesis says, and that the effect disappears along directions chosen at random or orthogonally.

\paragraph{7. Localising is not understanding.} \citet{meng2022locating} located factual associations in the
MLPs of middle layers by patching, and edited the facts by changing weights there. \citet{hase2023does} then
found that the layer indicated by patching does not predict which layer is best to edit. Knowing where an
intervention has an effect, and knowing how the computation works, are different things. \emph{Control:} do not
use the success of an edit as evidence for a mechanism, or the reverse.

\paragraph{8. A narrow distribution and a good story.} The IOI circuit was established on templated
sentences, and an attribution graph on one prompt. Labels such as ``name mover'' are summaries made by people,
and a label can be more general or more specific than what the component does. \emph{Control:} state the
distribution on which the claim was tested, derive a prediction for inputs outside it, test that prediction,
and report the share of the behaviour that the explanation leaves open.

\begin{exercise}[\extitle{Base rates}{derivation}]\label{ex:baserate}
A latent fires on $90\%$ of the texts about a certain concept and on $5\%$ of all other texts. The concept occurs
in $1\%$ of the corpus. (a) Which fraction of the texts on which the latent fires is about the concept?
(b) Suppose that the latent fires more strongly on texts about the concept than on other texts. What will the
list of its $20$ strongest activations look like, and what will a reader of that list conclude? (c) Which control
of failure mode 3 reveals the problem?
\end{exercise}

\subsection{Matching the evidence to the claim}

A useful habit is to ask of every claim, including one's own, which kind of evidence it would need.
\Cref{tab:evidence} gives the minimum for the four kinds of claims that occur most often.

\begin{table}[t]
\centering\small
\begin{tabular}{@{}p{0.34\textwidth}p{0.6\textwidth}@{}}
\toprule
Claim & Minimal evidence \\
\midrule
``Property $y$ is represented in activation $x$.'' & A probe or latent with held-out accuracy, a control task or
shuffled labels, and the same analysis on a model with random weights. \\ \addlinespace
``The model uses this representation.'' & An intervention on the representation that changes the behaviour as
predicted, with a random direction of equal size as comparison. \\ \addlinespace
``Component $g$ matters for task $T$.'' & Patching or ablation with a stated replacement value and a stated
distribution of prompts. Direct and total effect. \\ \addlinespace
``This circuit is how the model does $T$.'' & Faithfulness, completeness and minimality. A confirmed prediction
on new inputs. The size of the unexplained remainder. \\
\bottomrule
\end{tabular}
\caption{The weakest evidence that supports each common type of claim. A weaker kind of evidence supports only
a weaker claim.}
\label{tab:evidence}
\end{table}

\begin{exercise}[\extitle{A lie detector}{concept}]\label{ex:critique}
A blog post reports: ``We trained a linear probe on layer 20 of a chat model and can predict with 99\% accuracy
whether the model is about to lie. We have found the model's lie detector.'' List three experiments that are
missing, and say for each which failure mode of this section it addresses.
\end{exercise}

\begin{exercise}[\extitle{Three abstracts}{concept}]\label{ex:abstracts}
For each of the following summaries, name the row of \cref{tab:evidence} to which the claim belongs and the most
important piece of evidence that is missing. (a) ``Zero-ablating the MLP of layer 14 destroys the ability of the
model to add two-digit numbers. Addition is therefore computed in this MLP.'' (b) ``A latent of our SAE fires on
$85\%$ of the answers in which the model deceives the user. The model therefore knows when it is deceiving.''
(c) ``Our automatically discovered circuit reaches $98\%$ faithfulness under mean ablation. It contains $70\%$ of
all edges of the model.''
\end{exercise}

\begin{exercise}[\extitle{A baseline for steering}{concept}]\label{ex:steerbaseline}
Adding a steering vector raises the rate of a certain behaviour from $10\%$ to $60\%$. Adding a random vector of
the same norm raises it to $35\%$. (a) Which part of the increase can be credited to the direction? (b) Suggest a
reason why a random vector changes the behaviour at all. (c) Which further measurements are needed before the
vector is used to control the model?
\end{exercise}

\begin{takeaways}
Before believing an interpretability result, ask six questions. What exactly is the hypothesis, and what would
count against it? On which distribution of inputs was it tested? What was the counterfactual in each
intervention? Does the result disappear for a random model, random labels or a random direction? Is the
evidence observational or interventional, and does the claim go beyond it? How much of the behaviour is left
unexplained?
\end{takeaways}

\begin{readinglist}
\item \citet{adebayo2018sanity}, \emph{NeurIPS}, and \citet{heap2025sparse}, \emph{ICLR}. The same control, a model
with random weights, applied to attribution maps and to sparse autoencoders (failure mode 2).
\item \citet{hewitt2019designing}, \emph{EMNLP}. Control tasks, the standard guard against failure mode 1.
\item \citet{makelov2024subspace}, \emph{ICLR}. A complete worked example of failure mode 6, on the IOI task.
\item \citet{hase2023does}, \emph{NeurIPS}. The evidence behind failure mode 7.
\item \citet{friedman2024interpretability}, \emph{ICML}. Simplified descriptions of a model can agree with it on
the training distribution and disagree elsewhere (failure mode 8).
\item \citet{miller2024transformer}, \emph{COLM}. Shows how much reported faithfulness depends on details of the
ablation.
\item \citet{shi2024hypothesis}, \emph{NeurIPS}. Turns the criteria for circuits into statistical hypothesis
tests and applies them to published circuits.
\item \citet{meloux2025everything}, \emph{ICLR}. Enumerates all candidate explanations in small networks and finds
that many different ones pass the usual criteria. An explanation that passes need not be the only one.
\end{readinglist}
